% Syntra LaTeX Template
% For Arabic support, use XeLaTeX with:
%   \usepackage{fontspec}
%   \usepackage{polyglossia}
%   \setdefaultlanguage{english}
%   \setotherlanguage{arabic}
%   \newfontfamily\arabicfont[Script=Arabic]{Amiri}

\documentclass[12pt,a4paper]{article}

% Font and encoding
\usepackage[utf8]{inputenc}
\usepackage[T1]{fontenc}

% Math packages
\usepackage{amsmath,amssymb}

% Table support
\usepackage{longtable,booktabs}

% Graphics and links
\usepackage{graphicx}
\usepackage{hyperref}
\usepackage{geometry}

% Page geometry
\geometry{
    left=2.5cm,
    right=2.5cm,
    top=3cm,
    bottom=3cm
}

% Fix pandoc tightlist issue
\providecommand{\tightlist}{\setlength{\itemsep}{0pt}\setlength{\parskip}{0pt}}

% Hyperref settings
\hypersetup{
    colorlinks=true,
    linkcolor=blue,
    filecolor=magenta,
    urlcolor=cyan,
    pdftitle={Syntra Research Documentation},
    pdfauthor={Syntra-Env},
    pdfpagemode=FullScreen
}

\begin{document}

\section{Technical Implementation: The Conformance-First
Workflow}\label{technical-implementation-the-conformance-first-workflow}

To implement PRISM/UOR, you are not just writing tests; you are building
a \textbf{programmatic validation engine} that gates your repository.

\subsection{1. The Modeling Layer (The
Data)}\label{the-modeling-layer-the-data}

Do not just write functions. Define your ontology (data model) formally.
* \textbf{How:} Serialize your domain data into formats like Turtle
(\texttt{.ttl}) or JSON-LD. * \textbf{Where:} Store these in a dedicated
directory (e.g., \texttt{conformance/standards/} or \texttt{spec/}). *
\textbf{Mechanism:} Your application code should be derived from this
data.

\subsection{2. The Validation Layer (The
Engine)}\label{the-validation-layer-the-engine}

Do not rely on ad-hoc unit tests. Build a dedicated conformance
validator. * \textbf{How:} Write a tool (in Rust, or your language of
choice) that loads your domain data and runs formal mathematical
assertions against it. * \textbf{Mechanism:}
\texttt{bash\ \ \ \ \ \#\ Run\ the\ validator\ specifically\ \ \ \ \ cargo\ run\ -\/-bin\ uor-conformance}
* \textbf{Where:} Create a dedicated crate or module (e.g.,
\texttt{conformance/}) that acts as the ``source of truth'' for what is
compliant.

\subsection{3. The Generation \& Drift Layer (The
Sync)}\label{the-generation-drift-layer-the-sync}

If you generate code from your ontology, you must enforce
synchronization. * \textbf{How:} Create a generator (e.g.,
\texttt{uor-crate}) that outputs source code from your ontology
definitions. * \textbf{Mechanism:} In your CI, force a comparison
between the generated output and your checked-in code.
\texttt{bash\ \ \ \ \ \#\ Generate\ code\ \ \ \ \ cargo\ run\ -\/-bin\ uor-crate\ \ \ \ \ \#\ Fail\ if\ the\ generated\ code\ doesn\textquotesingle{}t\ match\ the\ repo\ \ \ \ \ git\ diff\ -\/-exit-code\ path/to/generated/code/}

\subsection{4. The Enforcement Layer (The
Gate)}\label{the-enforcement-layer-the-gate}

Configure your CI (\texttt{.github/workflows/ci.yml}) to be the final
authority. * \textbf{How:} Add your validation and drift checks as
required steps. * \textbf{Mechanism:}
\texttt{yaml\ \ \ \ \ steps:\ \ \ \ \ \ \ -\ name:\ Run\ conformance\ suite\ \ \ \ \ \ \ \ \ run:\ cargo\ run\ -\/-bin\ uor-conformance\ \ \ \ \ \ \ -\ name:\ Verify\ no\ drift\ in\ generated\ code\ \ \ \ \ \ \ \ \ run:\ git\ diff\ -\/-exit-code\ foundation/src/}
* \textbf{Result:} A Pull Request cannot be merged unless these steps
return an exit code of \texttt{0}.

\begin{center}\rule{0.5\linewidth}{0.5pt}\end{center}

\subsection{Summary for the Community}\label{summary-for-the-community}

{\def\LTcaptype{none} % do not increment counter
\begin{longtable}[]{@{}
  >{\raggedright\arraybackslash}p{(\linewidth - 2\tabcolsep) * \real{0.5000}}
  >{\raggedright\arraybackslash}p{(\linewidth - 2\tabcolsep) * \real{0.5000}}@{}}
\toprule\noalign{}
\begin{minipage}[b]{\linewidth}\raggedright
Goal
\end{minipage} & \begin{minipage}[b]{\linewidth}\raggedright
Technique
\end{minipage} \\
\midrule\noalign{}
\endhead
\bottomrule\noalign{}
\endlastfoot
\textbf{Model} & Define ontology in formal data (Turtle/JSON-LD). \\
\textbf{Validate} & Build a validator tool to assert axioms against that
data. \\
\textbf{Sync} & Generate code from ontology; fail CI if generated code
!= repo code. \\
\textbf{Gate} & Add these steps to \texttt{ci.yml}. \textbf{If they
fail, release is blocked.} \\
\end{longtable}
}

\textbf{The result:} You stop asking \emph{if} the code is correct. The
pipeline proves it is.

\end{document}
